\paragraph{Unnumbered continuation: actual filled carriers and finite controlled expansions.}
Retain the original data and pointed group $\mathscr F$ of
(CFT.1). Let $P$ be a compact metric space with specified point
$u_0$, and let $u\mapsto h_u\in\mathscr F$ be continuous in
the original closed-disk map-and-inverse topology, with
$h_{u_0}=\id$. Every map below uses this given family and its
specified lifts.

Fix a homeomorphism $\iota:\overline D\to\overline{\mathbb D}$,
and write
\[
 e(z,w)=|\iota z-\iota w|,\qquad
 e_Q(z,w)=\sup_{q\in Q}e(qz,qw).
                                                               \tag{FCE.1}
\]
This is a fixed coordinate on the original closed disk. It is
not a change of the operator action. Define the two comparison
functions for every $a\geq0$ by
\[
 \omega_{d,e}(a)=\max_{d(z,w)\leq a}e(z,w),\qquad
 \omega_{e,d}(a)=\max_{e(z,w)\leq a}d(z,w).
                                                               \tag{FCE.2}
\]
The maxima exist by compactness. Both functions are nondecreasing,
vanish at zero and tend to zero at zero: a contrary sequence has
a convergent subsequence of pairs with equal limit points. The
identity map and its inverse therefore obey
$e_Q\leq\omega_{d,e}\circ d_Q$ and
$d_Q\leq\omega_{e,d}\circ e_Q$. For the sums $D_e,D_d$ of
the original uniform map and inverse distances, the same
pointwise estimates give
$D_e\leq2\omega_{d,e}(D_d)$ and
$D_d\leq2\omega_{e,d}(D_e)$. These are both topology
comparisons, on the original maps with the original labels.

Put
\[
 a^+(y)=\max_{u\in P}e_Q(h_uz,z),\qquad
 a^-(y)=\max_{u\in P}e_Q(h_u^{-1}z,z),\quad y=\pi z.
                                                               \tag{FCE.3}
\]
Equivariance makes these independent of the representative $z$.
The proof of (CFT.12)--(CFT.15), applied to $e$ using (FCE.2),
shows that $a^+,a^-\in C_0(N)$. In particular both of the
actual displacements have a common continuous positive
$C_0$ majorant, for example $a^++a^-+\tau$, where $\tau$
is the explicit function constructed in (CFT.14).

\emph{Filling an actual compact connected set.}
For a nonempty compact connected $A\subset O$, let
$\operatorname{Fill}(A)$ be $A$ together with all components
of $O\setminus A$ whose closures in $O$ are compact.
In the coordinate $\iota$, choose $r<1$ with
$\iota A\subset\{|z|<r\}$. The connected annulus
$r\leq|z|<1$ lies in a single component $U_A$ of
$O\setminus A$. Every other component has closure contained
in the compact disk $|z|\leq r$; hence
$\operatorname{Fill}(A)=O\setminus U_A$ is compact.
It is connected because each filled component has boundary
in $A$, and that boundary is nonempty. For the nonempty-boundary
assertion, a component with empty boundary would be open and
closed in the connected space $O$ while missing the nonempty
set $A$.

Every point outside the Euclidean convex hull of $\iota A$
has a separating straight line and a ray to the unit circle
which misses that convex hull. It therefore lies in $U_A$.
Consequently
\[
 \iota\operatorname{Fill}(A)\subset
       \operatorname{conv}(\iota A),\qquad
 \operatorname{diam}_e\operatorname{Fill}(A)
       =\operatorname{diam}_e A.
                                                               \tag{FCE.4}
\]
To verify the diameter equality, two finite convex combinations
$\sum_i\alpha_i v_i$ and $\sum_j\beta_j w_j$ have difference
$\sum_{i,j}\alpha_i\beta_j(v_i-w_j)$, of norm at most
$\operatorname{diam}(\iota A)$. Thus the convex hull has this
same diameter; the reverse inequality uses $A\subset
\operatorname{Fill}(A)$.

A homeomorphism $q:O\to O$ maps complement components
bijectively and maps compact closures to compact closures;
its inverse proves the reverse implication. Thus the exact
set map and its inverse give
\[
 q\operatorname{Fill}(A)=\operatorname{Fill}(qA),\qquad
 \operatorname{diam}_{e_Q}\operatorname{Fill}(A)
       =\sup_{q\in Q}\operatorname{diam}_e(qA)
       =\operatorname{diam}_{e_Q}A.
                                                               \tag{FCE.5}
\]
This uses no extension of $q$ to the circle.

Let $A$ now be a compact arc with endpoint $w$ fixed by
every $h_u$. The compact connected set
$A_P=A\cup\bigcup_{u\in P}h_u(A)$ has that common point.
Joint continuity on $P\times A$ proves compactness even
when $P$ is not connected. Equivariance gives the exact formula
\[
 q\operatorname{Fill}(A_P)=
 \operatorname{Fill}\left(qA\cup\bigcup_{u\in P}h_u(qA)\right),
 \qquad
 \operatorname{diam}_{e_Q}\operatorname{Fill}(A_P)
 \leq\operatorname{diam}_{e_Q}A
            +2\sup_{y\in\pi A}a^+(y).
                                                               \tag{FCE.6}
\]
Indeed any two points of the union on the right are images of
two points of $qA$, or are themselves such points. The triangle
inequality through those two original points gives the displayed
bound, uniformly in $q$, and (FCE.4)--(FCE.5) transfer it to
the filled set. In particular the estimate concerns every
actual translated carrier in the original disk.

\emph{An explicit local expansion of a common fixed region.}
Choose an embedding
$\zeta:[-1,1]^2\to O$ whose closed image $D$ has pairwise
disjoint $Q$-translates and with $x\notin QD$.
Such disks exist in any sufficiently small covering chart
away from $Qx$. Let $K\subset O$ be closed and $Q$-invariant,
fixed pointwise by every $h_u$, and suppose its intersection
with the chosen disk is exactly
$\zeta\{(s,v):v\leq0\}$. These are the local fixed-set data
for this operation; no such chart for an arbitrary fixed set
is inferred.

Choose $0<a_0<1$ and a continuous
$\theta:[-1,1]\to[0,1]$ equal to zero near the two endpoints.
For $0\leq a\leq a_0$ set
\[
 F_a(v)=
 \begin{cases}a+(1+a)v,&-1\leq v\leq0,\\
               a+(1-a)v,&0\leq v\leq1,
 \end{cases}
 \quad
 F_a^{-1}(w)=
 \begin{cases}(w-a)/(1+a),&-1\leq w\leq a,\\
               (w-a)/(1-a),&a\leq w\leq1.
 \end{cases}                                                   \tag{FCE.7}
\]
Both composites are the identity on their two stated pieces.
The formulas agree at their breakpoints, and $1-a\geq1-a_0>0$.
Thus maps and inverses vary jointly continuously with $a$.
For $t\in[0,1]$ define
\[
 \tau_t\bigl(q\zeta(s,v)\bigr)
       =q\zeta\bigl(s,F_{t a_0\theta(s)}(v)\bigr),
       \quad q\in Q,
 \qquad
 \tau_t(z)=z\quad(z\notin QD).
                                                               \tag{FCE.8}
\]
The disjoint translates make this definition unambiguous.
It fixes every disk boundary; properness of $Q$ makes the
translates of the compact disk locally finite in $O$. The
inverse is the same formula with $F_a^{-1}$, pasted along
the same boundaries. These facts prove joint interior
continuity of both maps and inverses. Their quotient support
lies in the one compact set $\pi D$, so C.81 gives joint
extension by the identity to the original circle. They
commute with every original operator and fix $x$.

The exact image of $K$ in $qD$ is now
\[
 \tau_tK\cap qD
   =q\zeta\{(s,v):v\leq t a_0\theta(s)\},\qquad
 K\subset\tau_tK,
 \qquad
 H_{u,t}=\tau_t h_u\tau_t^{-1}.
                                                               \tag{FCE.9}
\]
Outside $QD$, $\tau_tK$ equals $K$. The inclusion follows
from $t a_0\theta(s)\geq0$. Each $H_{u,t}$ fixes
$\tau_tK$ pointwise because
$H_{u,t}(\tau_tk)=\tau_th_uk=\tau_tk$.
It lies in $\mathscr F$, has inverse
$\tau_t h_u^{-1}\tau_t^{-1}$, starts at $h_u$, and is the
identity for $u=u_0$. Thus it is an actual based deformation
which enlarges this common fixed region and retains $K$.

Here the topology is the original one. On a compact disk,
uniform convergence $f_n\to f$, $g_n\to g$ gives uniform
convergence of $f_ng_n$ to $fg$: the difference is bounded
by $\sup e(f_n,f)$ plus the modulus of uniform continuity
of $f$ applied to $\sup e(g_n,g)$. Apply this also to
$g_n^{-1}f_n^{-1}$ for the inverse maps. Therefore composition
and inversion are continuous in the stated map-and-inverse
metric. This proves continuity of $(u,t)\mapsto H_{u,t}$;
(FCE.2) gives both comparisons with the original metric $d$.

\emph{Two actual disk families and quantitative bounds.}
The same estimates apply to every jointly continuous
boundary-fixed isotopy on $D$ with its equivariant paste
$\tau_t$, whether or not it is the expansion (FCE.8).
Write $S=QD$ and
\[
 s_D=\operatorname{diam}_{e_Q}D,\quad
 m_u^+=\operatorname{diam}_{e_Q}h_u(D),\quad
 m_u^-=\operatorname{diam}_{e_Q}h_u^{-1}(D),\quad
 \alpha_D^\pm=\sup_{y\in\pi D}a^\pm(y).
                                                               \tag{FCE.10}
\]
The triangle inequality through two points of $qD$ gives
$m_u^\pm\leq s_D+2\alpha_D^\pm$ for every $u$.
The exact left correction and its inverse are
\[
 A_{u,t}=H_{u,t}h_u^{-1}
       =\tau_t\,(h_u\tau_t^{-1}h_u^{-1}),\qquad
 A_{u,t}^{-1}=h_u\tau_t h_u^{-1}\tau_t^{-1}.
                                                               \tag{FCE.11}
\]
The first factor is pasted on the disks $qD$.
The second is pasted on the disks
$h_u(qD)=qh_u(D)$, which are pairwise disjoint because
$h_u$ is injective. The individual maps are exactly the
displayed conjugates; the two families may overlap with
one another. This factorization retains all original deck
labels without choosing any new point tracks.

Every displacement of $\tau_t$ or $\tau_t^{-1}$ in the
metric $e$ is at most $s_D$. Every displacement of
$h_u\tau_t^{\pm1}h_u^{-1}$ is at most $m_u^+$.
The triangle inequality for the two factors in (FCE.11),
and then the same argument with $h_u^{-1}$ in place of
$h_u$, proves
\[
 \begin{aligned}
 \sup_{z\in\overline D}e(H_{u,t}z,h_uz)&\leq s_D+m_u^+,\\
 \sup_{z\in\overline D}e(H_{u,t}^{-1}z,h_u^{-1}z)&\leq s_D+m_u^-,\\
 D_e(H_{u,t},h_u)&\leq2s_D+m_u^++m_u^-\\
                 &\leq4s_D+2\alpha_D^++2\alpha_D^-.
 \end{aligned}                                                 \tag{FCE.12}
\]
In the original $d$ metric the first two bounds instead give
$D_d(H_{u,t},h_u)\leq
\omega_{e,d}(s_D+m_u^+)+\omega_{e,d}(s_D+m_u^-)$.
Thus both maps and inverses are controlled by the actual
common displacement envelope on this carrier.

The exact input regions on which a change can occur are
\[
 \begin{aligned}
 \{z:H_{u,t}z\ne h_uz\}&\subset S\cup h_u^{-1}S,\\
 \{z:H_{u,t}^{-1}z\ne h_u^{-1}z\}&\subset S\cup h_uS,\\
 \operatorname{supp}_O(A_{u,t})&\subset S\cup h_uS.
 \end{aligned}                                                 \tag{FCE.13}
\]
For example outside $S\cup h_u^{-1}S$ the first
$\tau_t^{-1}$ fixes the input and the last $\tau_t$
fixes its $h_u$ image. The inverse formula proves the
second inclusion; the two factors prove the third.
Here $S$ is closed in $O$ by local finiteness of the
closed disks, and $h_uS$ is closed by the homeomorphism
$h_u$. Thus taking closure in $O$ preserves the third
inclusion.
All these quotient carriers lie in compact sets uniformly
over $u\in P$, because $\pi D$ is compact and the given
quotient maps and inverses evaluate continuously on it.

\emph{An actual small-displacement family requiring more than one such disk paste.}
For this example take the original disk to be the closed
unit disk and $d=e$ to be Euclidean distance. Use the exact
radial map
$j:\mathbb R^2\to\mathbb D$, $j(z)=z/(1+|z|)$, with
inverse $j^{-1}(w)=w/(1-|w|)$, and let
$q=jT_{(1,0)}j^{-1}$. The group $Q=\langle q\rangle$
is a free proper closed discrete action; conjugation by
$j$ verifies each property from integer translations.
All plane translations conjugated by $j$ extend by the
identity on the original circle, commute with $q$, and
act transitively. To check that extension directly, a
bounded translation vector changes $j(z)$ by a quantity
tending uniformly to zero as $|z|\to\infty$, by the
estimate below. Thus this is an actual C.81 instance.

Put $\beta(v)=\max(0,1-4|v|)$. For $n\geq1$ and
$0\leq s\leq1$, the actual maps and inverses are
\[
 F_{n,s}(a,y)=(a+s\beta(y-n),y),\qquad
 F_{n,s}^{-1}(a,y)=(a-s\beta(y-n),y),\qquad
 H_{n,s}=jF_{n,s}j^{-1}.
                                                               \tag{FCE.14}
\]
They commute with $q$ and fix $x=j(0,0)$. They extend by
the identity to the circle. Indeed, direct subtraction
gives, for $v,w\in\mathbb R^2$,
$|j(v)-j(w)|\leq
2|v-w|/(1+\min(|v|,|w|))$.
On the support of (FCE.14) both points have
$|y|\geq n-1/4$ and their horizontal difference is at most $1$.
The estimate applies equally to the inverse and proves
\[
 \sup e(H_{n,s}z,z)\leq\frac{2}{n+3/4},\qquad
 D_e(H_{n,s},\id)\leq\frac{4}{n+3/4}.
                                                               \tag{FCE.15}
\]
Under $\pi j(a,y)=(a\bmod\mathbb Z,y)$ the support is
in the compact strip
$(\mathbb R/\mathbb Z)\times[n-1/4,n+1/4]$.

Let $I_n=[2^{-n-1},2^{-n}]$, and let $s_n$ be the
piecewise linear function on $I_n$ which is zero at its
two endpoints and one at its midpoint $t_n$.
On $[0,1]$ set $H_t=H_{n,s_n(t)}$ for $t\in I_n$,
and $H_t=\id$ elsewhere. Both formulas agree at every
shared endpoint. Formula (FCE.15) proves continuity at
the accumulation point $0$ in the original map-and-inverse
metric. Identifying $0$ and $1$ gives an actual based loop in
$\mathscr F$. The explicit based nullhomotopy replaces
$s_n(t)$ by $(1-r)s_n(t)$, $0\leq r\leq1$; the same
uniform estimate proves joint continuity at $0$.
The family has a common $C_0$ displacement tail: outside
$(\mathbb R/\mathbb Z)\times[-1,N+1]$ only strips with
$n>N$ contribute, and their bound is
$2/(N+7/4)$, for maps and inverses alike.

At the midpoint, for $v_n=j(0,n)$ one has the literal
deck-step equality $H_{t_n}(v_n)=qv_n\ne v_n$.
A closed Jordan disk containing both points in its
interior meets its $q$-translate at $qv_n$.
More generally, a map pasted into a $Q$-invariant family
of pairwise disjoint disk interiors and fixing every
disk boundary preserves each disk. If its left
composition with $H_{t_n}$ were to fix $v_n$, it would
have to send $qv_n$ to $v_n$ inside one such disk.
That disk and its $q$-translate would both contain
$qv_n$ in their interiors. They are distinct: otherwise
all $q^kv_n$ would lie in one compact disk in $O$,
contrary to properness of this integer-translation action.
This contradicts the stated disjointness. The corresponding right correction would
send $v_n$ to $q^{-1}v_n$ and has the same contradiction.
Thus one disk-paste stage of this specified form cannot
make that correction, even for this continuous based
family with the explicit nullhomotopy above.

Formulas (FCE.11)--(FCE.13) provide two separately
disjoint actual disk families and allow overlap between
the stages. Formulas (FCE.4)--(FCE.6) control the filled
size of actual arcs and their given images, but do not
assert pairwise disjointness of those filled sets.
This note proves the finite operations and their exact
original-topology bounds. It constructs no infinite
sequence of compatible fixed-set expansions across all
edges or essential cycles, and makes no full noncompact
contraction assertion. The exact infinite-evaluation
obstruction in C.83 is unchanged.
